% Evaluation methodology. Resource protocol as implemented in paper/eval/run_efficiency.py and ann_scaling.py.
\section{Evaluation methodology}\label{sec:method}

\subsection{Systems and units}
All systems use the same archive-index search implementation with explicit configuration switches. The comparison includes plain and lemmatized BM25, full-dimensional EmbeddingGemma~2, a 256-dimensional EG2 variant, multilingual-E5, reciprocal-rank fusion, the adaptive system and three ablations removing phrase reranking, identifier handling or short-fragment filtering. Runs are document-level: the first retrieved chunk determines each document's rank, and subsequent chunks of that document are collapsed. Up to 100 documents are retained per query. The E5 index is text-only, unlike the full EG2 indexes; the differing candidate universes are reported as a limitation.

\subsection{Query population and metrics}
Retrieval means use the same 170 answerable queries for every system and judge variant, with each query weighted equally. The 30 designated no-answer queries are retained in the relevance files but excluded from those means; empty retrieval and positive-labelled results on them are reported separately. An answerable query with no positive pooled judgement remains in the average with score zero.

The primary metric is nDCG@10, using linear relevance gains 0, 1 and 2, discounted by $\log_2(i+1)$ at rank $i$ and normalized by the ideal ranking of the judged documents. Secondary metrics are reciprocal rank at 10, success at 1 and 5, and recall at 100. Binary metrics treat grades 1 and 2 as relevant. Every system's retrieved top-10 documents are judged, but judgements outside the depth-10 union are incomplete. Recall at 100 is therefore relative to the known relevant pool, not exhaustive corpus recall. Results are also grouped by query category and language.

\subsection{Uncertainty and comparisons}
Confidence intervals for query means and paired mean differences use a percentile bootstrap over queries with 10\,000 resamples and seed 2026. Two-sided paired sign-flip randomization tests use 10\,000 resamples and the add-one correction $(b+1)/(10\,000+1)$. Each metric and judge variant has a family of nine comparisons against the full adaptive system; p-values receive Holm adjustment within that family, with threshold 0.05. The original morphology research question is additionally assessed by a separately labelled paired lemma-versus-plain BM25 contrast on nDCG@10; its unadjusted p-value is not part of the adaptive comparison family. Neither separate confidence-interval overlap nor failure to reject a difference is interpreted as equivalence.

The initial inter-judge agreement report uses linear-weighted Cohen's kappa with a 95\% interval from 2\,000 query-level bootstrap resamples. Resampling queries, rather than individual judgement pairs, preserves within-query clustering. These agreement intervals are separate from the 10\,000-resample retrieval intervals. Rankings and pairwise ordering changes are reported for each relevance variant; numerically equal means within $10^{-12}$ share a rank.

\subsection{Reproducibility and evidence}
Initial judgement sheets, query definitions, corpus manifest, retrieval runs and available call logs are frozen before the reported aggregate effectiveness comparison is computed. A SHA-256 manifest validates those inputs; a separate prepared-review record binds the review to its original disagreements and presentation order. Structured call responses are logged with model, conversation and usage metadata when available. Canonical labels replace only disagreement grades after complete review, preserving every agreed grade. Report tables and figures are generated from the resulting JSON artifacts rather than maintained as independent numerical copies.

Unit and integration checks establish implementation behaviour; the public benchmark provides model-labelled retrieval evidence; the separate multimodal experiments use source-dataset labels. These evidence types do not substitute for a human relevance audit or deployment measurements.

\subsection{Resource measurements}
Query latency is the wall time of query embedding and chunk search in the evaluation harness, including lexical and semantic retrieval and fusion with a candidate limit of 400 chunks. The timer stops before document collapse and path resolution; these measurements are not end-to-end interface latency. Each system answers all 200 queries once as a warm-up and then once more under measurement, in one process, first with the GPU visible and then with \code{CUDA\_VISIBLE\_DEVICES} empty (the embedding models then run on the CPU in float32). During each measured pass the harness records per-core CPU use from \code{/proc/stat}, CPU seconds per thread of the process from \code{/proc/self/task}, the change in interrupt counts per core from \code{/proc/interrupts}, the peak resident set size, and GPU memory sampled every 0.2~s with \code{nvidia-smi}. The machine was an ordinary desktop session; the one-minute load average before each step is written to the measurement log.

Indexing cost is measured on a fixed random sample of the corpus (seed 2026: 300 notes, 40 documents, 15 tables, 60 pictures and 60 recordings) indexed by the \code{archive} command into an empty database with the GPU. The CPU run uses a fifth-sized sample: 60 notes, 8 documents, 3 tables, 12 pictures and 12 recordings. The two samples differ in size and composition, so their wall times do not estimate a paired CPU/GPU speedup. The sizes of the four full benchmark indexes are reported as built.

For quantization, the text chunks of the full index (84\,788 vectorized rows in text collections) are embedded again by GGUF builds of EmbeddingGemma~2 (f16, Q8\_0, Q4\_K\_M) served by llama.cpp with mean pooling, with the same \code{title: \ldots{} | text: \ldots} prefixes; queries use the same retrieval prompt. Dense retrieval with each build runs on the 4\,137 text documents and is compared with the SentenceTransformers bfloat16 model on the same documents (\sys{dense\_eg2\_text}) under the same labels. Documents that a quantized run brings into its top~10 but that no judge saw are counted as non-relevant, and their share is reported. The text collections include 21 PDF-page rows embedded as images by the full model. GGUF instead encodes their stored text (485 characters in total); 84\,767 rows are text inputs. Thus identical document candidates do not mean identical input modality for every chunk, and the comparison does not isolate quantization alone.

Scaling of exact dense search is measured outside archive-index with collections grown from the 86\,726 real EG2 vectors to $10^5$, $3\cdot10^5$, $10^6$ and $2\cdot10^6$ vectors; the added vectors are noisy copies of real ones (Gaussian noise, $\sigma=0.05$, renormalized), so the collections are synthetic. Each of the 200 UA-PAR query vectors is searched one at a time with an optimized exact NumPy top-10 selection (matrix product and \code{argpartition}), an exact Faiss index and an HNSW index ($M=32$, construction depth 200, search depth 64 and 256). A sequential follow-up calls the actual archive-index vector-ranking method on the same matrices: candidate gathering, full \code{argsort} and construction of 400 result tuples, excluding the SQLite cache-version poll. This distinguishes the optimized reference from the deployed algorithm. HNSW recall@10 counts matching vector IDs against the NumPy reference; ties and floating-point differences are not credited as equivalent IDs. Agreement between the exact Faiss and NumPy rankings is audited separately, together with the fraction of Faiss results at or above the reference tenth score within $10^{-6}$.

\subsection{Software and hardware}
All measurements ran on one laptop: Intel Core i9-14900HX (32 logical CPUs), 31~GiB RAM, NVIDIA GeForce RTX 4060 Laptop GPU with 8~GB, Linux 7.0. Software: Python 3.13.14, PyTorch 2.14.1, Transformers 5.19.0, SentenceTransformers 6.1.0, NumPy 2.5.4, SQLite 3.53.1, pymorphy3 2.0.6, ir-measures 0.4.3, llama.cpp commit \code{8bb3cf0}, and Faiss 1.15.1 in a separate environment for the scaling experiment. The full list is stored in \code{efficiency.json}.
